\capitulofrontal{O que cai: edital, matriz e raio-x}{matriz}

O ENEM é organizado pela \textbf{Matriz de Referência}, publicada pelo INEP e adotada nos editais do
exame. Para Matemática e suas Tecnologias, ela define \textbf{7 competências} e \textbf{30 habilidades}
e lista, em anexo, os \textbf{objetos de conhecimento} — os conteúdos que podem ser cobrados.
Esta coleção cobre \emph{todos} esses objetos, distribuídos em três volumes.

\secaofrontal{Objetos de conhecimento do edital e onde estudá-los}

Clique nos selos para ir ao capítulo. Os selos com \faExternalLink* abrem outro volume da coleção
(os três arquivos precisam estar na mesma pasta; alguns leitores de PDF pedem confirmação).

{\small
\renewcommand\arraystretch{1.32}
\begin{longtable}{@{}>{\raggedright\arraybackslash}p{0.55\linewidth}>{\raggedright\arraybackslash}p{0.41\linewidth}@{}}
\toprule
\textbf{Objeto de conhecimento (edital)} & \textbf{Onde estudar} \\
\midrule
\endhead
\multicolumn{2}{@{}l}{\textbf{\color{ebookPrim}Conhecimentos numéricos}}\\
Operações em conjuntos numéricos (naturais, inteiros, racionais e reais) & \caplink{1}{1} \\
Desigualdades & \caplink{1}{2} \\
Divisibilidade e fatoração & \caplink{1}{3} \\
Razões e proporções & \caplink{1}{4} \\
Porcentagem e juros & \caplink{1}{6}\ \caplink{1}{7} \\
Relações de dependência entre grandezas & \caplink{1}{5} \\
Sequências e progressões & \caplink{1}{8} \\
Princípios de contagem & \caplink{1}{9} \\[2pt]
\multicolumn{2}{@{}l}{\textbf{\color{ebookPrim}Conhecimentos geométricos}}\\
Características das figuras geométricas planas e espaciais & \caplink{2}{4}\ \caplink{2}{11} \\
Grandezas, unidades de medida e escalas & \caplink{2}{1}\ \caplink{2}{2} \\
Comprimentos, áreas e volumes & \caplink{2}{9}\ \caplink{2}{10}\ \caplink{2}{12} \\
Ângulos; posições de retas & \caplink{2}{3} \\
Simetrias de figuras planas ou espaciais & \caplink{2}{5} \\
Congruência e semelhança de triângulos; teorema de Tales & \caplink{2}{6} \\
Relações métricas nos triângulos & \caplink{2}{7} \\
Circunferências & \caplink{2}{9} \\
Trigonometria do ângulo agudo & \caplink{2}{8} \\[2pt]
\multicolumn{2}{@{}l}{\textbf{\color{ebookPrim}Conhecimentos de estatística e probabilidade}}\\
Representação e análise de dados & \caplink{1}{10} \\
Medidas de tendência central (médias, moda e mediana) & \caplink{1}{11} \\
Desvios e variância & \caplink{1}{12} \\
Noções de probabilidade & \caplink{1}{13} \\[2pt]
\multicolumn{2}{@{}l}{\textbf{\color{ebookPrim}Conhecimentos algébricos}}\\
Gráficos e funções & \caplink{3}{3} \\
Funções algébricas do 1.º e do 2.º graus & \caplink{3}{4}\ \caplink{3}{5} \\
Funções polinomiais e racionais & \caplink{3}{6} \\
Funções exponenciais e logarítmicas & \caplink{3}{7}\ \caplink{3}{8} \\
Equações e inequações & \caplink{3}{1} \\
Relações no ciclo trigonométrico e funções trigonométricas & \caplink{3}{9} \\[2pt]
\multicolumn{2}{@{}l}{\textbf{\color{ebookPrim}Conhecimentos algébricos/geométricos}}\\
Plano cartesiano; retas; paralelismo e perpendicularidade & \caplink{3}{10} \\
Circunferências & \caplink{3}{11} \\
Sistemas de equações & \caplink{3}{2} \\
\bottomrule
\end{longtable}
}

\secaofrontal{Raio-x deste volume}

Quantas questões oficiais (\anosbanco) tiveram como foco principal cada conteúdo deste volume, com a
divisão por nível de dificuldade TRI. Dos \totalbanco\ itens válidos do período, \textbf{\csname volstats@\numvolume\endcsname} pertencem a conteúdos deste volume.
Clique no nome para ir ao capítulo.

\begin{center}
\raioxvolume
\end{center}
{\scriptsize\centering
\textcolor{nivelA}{\faSquare}\,Fácil\quad\textcolor{nivelB}{\faSquare}\,Média\quad
\textcolor{nivelC}{\faSquare}\,Difícil\quad\textcolor{nivelD}{\faSquare}\,Muito difícil\par}

\secaofrontal{Competências e habilidades da Matriz de Referência}

As habilidades aparecem no cabeçalho de cada questão (por exemplo, \chip{ebookPrim}{H16}).
A barra mostra quantas questões oficiais de \anosbanco\ avaliaram cada habilidade.

{\small
\renewcommand\arraystretch{1.22}
\begin{longtable}{@{}p{7mm}>{\raggedright\arraybackslash}p{0.66\linewidth}l@{}}
\multicolumn{3}{@{}>{\raggedright\arraybackslash}p{\linewidth}@{}}{\textbf{\color{ebookDark}Competência 1} — Construir significados para os números naturais, inteiros, racionais e reais.}\\[2pt]
\habitem{1}{Reconhecer, no contexto social, diferentes significados e representações dos números e operações — naturais, inteiros, racionais ou reais.}
\habitem{2}{Identificar padrões numéricos ou princípios de contagem.}
\habitem{3}{Resolver situação-problema envolvendo conhecimentos numéricos.}
\habitem{4}{Avaliar a razoabilidade de um resultado numérico na construção de argumentos sobre afirmações quantitativas.}
\habitem{5}{Avaliar propostas de intervenção na realidade utilizando conhecimentos numéricos.}
\multicolumn{3}{@{}>{\raggedright\arraybackslash}p{\linewidth}@{}}{\textbf{\color{ebookDark}Competência 2} — Utilizar o conhecimento geométrico para realizar a leitura e a representação da realidade e agir sobre ela.}\\[2pt]
\habitem{6}{Interpretar a localização e a movimentação de pessoas/objetos no espaço tridimensional e sua representação no espaço bidimensional.}
\habitem{7}{Identificar características de figuras planas ou espaciais.}
\habitem{8}{Resolver situação-problema que envolva conhecimentos geométricos de espaço e forma.}
\habitem{9}{Utilizar conhecimentos geométricos de espaço e forma na seleção de argumentos propostos como solução de problemas do cotidiano.}
\multicolumn{3}{@{}>{\raggedright\arraybackslash}p{\linewidth}@{}}{\textbf{\color{ebookDark}Competência 3} — Construir noções de grandezas e medidas para a compreensão da realidade e a solução de problemas do cotidiano.}\\[2pt]
\habitem{10}{Identificar relações entre grandezas e unidades de medida.}
\habitem{11}{Utilizar a noção de escalas na leitura de representação de situação do cotidiano.}
\habitem{12}{Resolver situação-problema que envolva medidas de grandezas.}
\habitem{13}{Avaliar o resultado de uma medição na construção de um argumento consistente.}
\habitem{14}{Avaliar proposta de intervenção na realidade utilizando conhecimentos geométricos relacionados a grandezas e medidas.}
\multicolumn{3}{@{}>{\raggedright\arraybackslash}p{\linewidth}@{}}{\textbf{\color{ebookDark}Competência 4} — Construir noções de variação de grandezas para a compreensão da realidade e a solução de problemas do cotidiano.}\\[2pt]
\habitem{15}{Identificar a relação de dependência entre grandezas.}
\habitem{16}{Resolver situação-problema envolvendo a variação de grandezas, direta ou inversamente proporcionais.}
\habitem{17}{Analisar informações envolvendo a variação de grandezas como recurso para a construção de argumentação.}
\habitem{18}{Avaliar propostas de intervenção na realidade envolvendo variação de grandezas.}
\multicolumn{3}{@{}>{\raggedright\arraybackslash}p{\linewidth}@{}}{\textbf{\color{ebookDark}Competência 5} — Modelar e resolver problemas que envolvem variáveis socioeconômicas ou técnico-científicas, usando representações algébricas.}\\[2pt]
\habitem{19}{Identificar representações algébricas que expressem a relação entre grandezas.}
\habitem{20}{Interpretar gráfico cartesiano que represente relações entre grandezas.}
\habitem{21}{Resolver situação-problema cuja modelagem envolva conhecimentos algébricos.}
\habitem{22}{Utilizar conhecimentos algébricos/geométricos como recurso para a construção de argumentação.}
\habitem{23}{Avaliar propostas de intervenção na realidade utilizando conhecimentos algébricos.}
\multicolumn{3}{@{}>{\raggedright\arraybackslash}p{\linewidth}@{}}{\textbf{\color{ebookDark}Competência 6} — Interpretar informações de natureza científica e social obtidas da leitura de gráficos e tabelas, realizando previsão de tendência, extrapolação, interpolação e interpretação.}\\[2pt]
\habitem{24}{Utilizar informações expressas em gráficos ou tabelas para fazer inferências.}
\habitem{25}{Resolver problema com dados apresentados em tabelas ou gráficos.}
\habitem{26}{Analisar informações expressas em gráficos ou tabelas como recurso para a construção de argumentos.}
\multicolumn{3}{@{}>{\raggedright\arraybackslash}p{\linewidth}@{}}{\textbf{\color{ebookDark}Competência 7} — Compreender o caráter aleatório e não determinístico dos fenômenos naturais e sociais e utilizar instrumentos adequados para medidas, determinação de amostras e cálculos de probabilidade para interpretar informações de variáveis apresentadas em uma distribuição estatística.}\\[2pt]
\habitem{27}{Calcular medidas de tendência central ou de dispersão de um conjunto de dados expressos em uma tabela de frequências de dados agrupados (não em classes) ou em gráficos.}
\habitem{28}{Resolver situação-problema que envolva conhecimentos de estatística e probabilidade.}
\habitem{29}{Utilizar conhecimentos de estatística e probabilidade como recurso para a construção de argumentação.}
\habitem{30}{Avaliar propostas de intervenção na realidade utilizando conhecimentos de estatística e probabilidade.}
\end{longtable}
}
{\footnotesize\color{ebookMuted}Fonte: INEP. \emph{Matriz de Referência ENEM} — Matemática e suas
Tecnologias e Anexo (objetos de conhecimento associados).\par}
